\documentclass{situationist-preprint}

\title{Nostalgia and Misrepresentation}
\msid{MND-17-b}
\msfamily{MND-17}
\msdiscipline{Philosophy of Mind}
\mstype{Consequence}
\mskeywords{nostalgia; emotional content; misrepresentation; fittingness; recalcitrant emotion; memory}

\begin{document}
\maketitle

\begin{abstract}
If emotions represent evaluative properties, as MND-17-a argues, they can misrepresent. We derive the content of nostalgia and its conditions of misrepresentation. Nostalgia represents a part of one's own past as having contained a good that is no longer available. It misrepresents in two ways, by representing a good that was not there and by representing as lost a good that is still available. We distinguish both from excess, which is a failure of proportion rather than of content, and we show that recalcitrant nostalgia, which persists against the belief that the past was bad, confirms that its content is not that of a belief. Nostalgia is thus correct when it tracks a real past good that is gone, which is the condition under which EXI-14-a holds it rational.
\end{abstract}

\section{The content of nostalgia}

MND-17-a argues that each emotion type represents its object as having the evaluative property that is its formal object (Theorem~3.1), on the ground that emotions are assessable for fit by a type-specific property (Principle~P1), that this assessment is not that of an accompanying belief (Principle~P2) or of a bodily feeling (Principle~P3), and that a state whose correctness depends on its object's having a property represents it as having it (Assumption~2.3). We apply the account to nostalgia.

A nostalgic episode is directed at a part of the subject's past: a childhood house, a period of friendship, a city as it was. It is assessed for fit in characteristic terms. We say that someone's nostalgia for her school years is misplaced when the years were miserable, and that nostalgia for a friendship is premature when the friend is still there to be visited. The two terms of the assessment are the past good and its loss. The bittersweet character of the emotion corresponds to them: the sweetness to the good, the bitterness to its absence \citep{Sedikides2008}.

\begin{proposition}\label{p:content}
Nostalgia represents a part of the subject's own past as having contained a good, and that good as no longer available to her.
\end{proposition}

\begin{proof}
By Principle~P1 of MND-17-a, fittingness for nostalgia depends on whether its object has the type-specific property, and the assessments just described show that property to be the conjunction of past goodness and present unavailability. By Assumption~2.3 of MND-17-a, nostalgia represents its object as having that property.
\end{proof}

Proposition~\ref{p:content} is Definition~2.1 of EXI-14-a, here derived from the general account of emotional content rather than stipulated. The formal object of nostalgia is thus the lost past good, as that of grief is the loss of a valued person and that of fear the dangerous.

\section{Two ways to misrepresent}

\begin{corollary}\label{c:mis}
A nostalgic episode misrepresents its object if and only if either (a) the remembered part of the past did not contain the good, or not as the episode represents it, or (b) the good is still available to the subject.
\end{corollary}

\begin{proof}
The content of Proposition~\ref{p:content} is a conjunction; the episode misrepresents just when a conjunct is false.
\end{proof}

Case (a) is the familiar complaint against nostalgia: it gilds the past. The psychological evidence for a \enquote{rosy view} of past events, on which events are evaluated more favorably in retrospect than they were while they happened, documents this source of error \citep{Mitchell1997}. Case (b) is less often noticed. Nostalgia for a place one could return to, or for a practice one could resume, represents a loss that has not occurred, and its correction is action rather than revision of memory.

Neither case is intrinsic to nostalgia. That an emotion type can misrepresent does not show that its instances usually do; fear misrepresents in phobia, yet fear is often fitting. Whether a nostalgic episode is correct is settled by the facts about the past and the subject's present access to it, as EXI-14-a holds: nostalgia is correct when it tracks a real good of the past that is no longer available (Lemma~3.1 there).

\begin{remark}\label{r:excess}
Misrepresentation must be distinguished from excess. An episode can be correct, the good real and gone, and still out of proportion to the good, or so absorbing that it crowds out response to present goods. Excess is a failure of proportion, not of content, and it is the second condition of rational nostalgia in Theorem~3.2 of EXI-14-a \citep{deSousa1987}.
\end{remark}

\section{Recalcitrant nostalgia}

A subject can feel nostalgia for a period she knows to have been bad: she sincerely judges that her years in a certain town were unhappy, and yet a song from that time moves her to longing for them. This is a recalcitrant emotion, in conflict with the subject's considered belief \citep{Brady2009}.

\begin{objection}[Recalcitrance shows no content]
If nostalgia can persist against the belief that the past was not good, it does not represent the past as good; it is a mood attached to memory, without correctness conditions.
\end{objection}

\begin{reply}
The conflict is evidence for content, not against it. The subject herself experiences the episode as at odds with her judgment, and she can be criticized for it, as one is for fear of a spider one knows to be harmless. There is a conflict only if the emotion presents things as being otherwise than the belief does. This is Principle~P2 of MND-17-a: the correctness conditions of an emotion are not those of an accompanying belief. Often, moreover, recalcitrant nostalgia is not wholly incorrect. The miserable years may have contained a smaller good, a friendship or a kind of freedom, that the episode picks out and the global judgment discounts; in that case the episode represents a real good under a misleading description of its extent, and the error is one of proportion.
\end{reply}

\begin{objection}[Nostalgia for what one never had]
People feel nostalgia for periods they did not live through. Such episodes cannot represent a part of the subject's own past.
\end{objection}

\begin{reply}
They are a different, if related, emotion: a longing for a past imagined as good, which represents an historical period rather than the subject's past as containing a good now unavailable. Its correctness conditions differ accordingly, and the term \enquote{nostalgia} applies to it by extension. Proposition~\ref{p:content} concerns the personal emotion, whose object is the subject's own past.
\end{reply}

\section{Conclusion}

From the thesis that emotions represent evaluative properties it follows that nostalgia represents a part of one's past as having contained a good that is no longer available, and that it misrepresents either by inventing or inflating the good or by mourning a good still at hand. Excess is a separate failure. Correct nostalgia tracks a real past good that is gone, and that is the condition under which nostalgia is rational.

\begin{thebibliography}{9}\small
\bibitem[Brady(2009)]{Brady2009} Brady, M.~S. (2009). The irrationality of recalcitrant emotions. \emph{Philosophical Studies}, 145(3), 413--430.
\bibitem[de Sousa(1987)]{deSousa1987} de Sousa, R. (1987). \emph{The Rationality of Emotion}. MIT Press.
\bibitem[Mitchell et~al.(1997)]{Mitchell1997} Mitchell, T.~R., Thompson, L., Peterson, E. and Cronk, R. (1997). Temporal adjustments in the evaluation of events: The \enquote{rosy view}. \emph{Journal of Experimental Social Psychology}, 33(4), 421--448.
\bibitem[Sedikides et~al.(2008)]{Sedikides2008} Sedikides, C., Wildschut, T., Arndt, J. and Routledge, C. (2008). Nostalgia: Past, present, and future. \emph{Current Directions in Psychological Science}, 17(5), 304--307.
\end{thebibliography}

\end{document}
